\section{Introduction}
\label{sec:introduction}

Sharding increases the capacity of distributed ledgers by partitioning state
and transaction processing across multiple shards. This parallelism creates a
coordination problem whenever a transaction spans more than one shard: the
system must preserve a consistent terminal outcome despite delayed or
duplicated messages, unmet quorum conditions, or interrupted execution.
Chainspace, OmniLedger, Prophet, CSLAP, and LightCross illustrate different
approaches to cross-shard transaction processing and atomic commit
\cite{albassam2018chainspace,kokoriskogias2018omniledger,hong2023prophet,du2024cslap,qi2024lightcross}.
Their correctness arguments, implementation evidence, and evaluation
methodologies differ substantially.

Formal methods provide evidence about protocol behavior beyond conventional
testing. TLA+ and related model-checking workflows analyze distributed
protocols and connect formal models with executable systems
\cite{schultz2025mongodb,cirstea2025traces,hackett2025tracelink,howard2025ccf,wang2023mocket,gao2025imocket}.
Mutation-based assessment instead asks whether a specification or property
suite reacts when selected controls are intentionally weakened
\cite{cordy2023mutation}. These techniques are complementary: bounded model
checking evaluates declared properties in finite configurations, mutation
analysis measures sensitivity to controlled defects, and trace conformance
checks selected implementation observations against a formal abstraction.

This paper studies an executable cross-shard commit protocol implemented in
DTL-Lab, a research-oriented Java simulator rather than a production
blockchain. The protocol coordinates source locking, receipt creation and
delivery, destination preparation, commit, abort, timeout, rollback, replay
protection, and quorum enforcement. Bounded TLA+ and Alloy representations
model the same protocol concerns through different state representations and
exploration semantics. The study assumes neither semantic equivalence between
the formal models nor intrinsic superiority of either tool.

The methodological objective is to connect several forms of evidence under one
frozen experimental protocol. Within the comparator set for which full text was
verified in our structured review, we did not identify a cross-shard study
reporting together bounded property verification with complementary TLA+ and
Alloy models, targeted mutation-based property validation, bounded
implementation-model trace conformance over valid and deliberately corrupted
traces, explicit verification-cost characterization, and independent
reproduction of the analytical artifact. This statement is limited to the
reviewed comparator set and is not a priority claim about the complete
literature. The contribution is therefore an integrated and reproducible
evidence workflow rather than novelty of any individual technique.

The study is organized around four research questions:

\begin{description}
\item[\textbf{RQ1.}] What bounded property-verification outcomes are obtained
for the valid TLA+ and Alloy models across the declared experimental
configurations?

\item[\textbf{RQ2.}] Do the evaluated TLA+ and Alloy properties detect
predefined scientific mutations that remove selected protocol controls?

\item[\textbf{RQ3.}] Do selected valid and deliberately corrupted Java
executions satisfy the expected bounded implementation-model trace-conformance
outcomes under the declared projection and scenario catalogue?

\item[\textbf{RQ4.}] How does verification cost vary across the evaluated
bound profiles and TLC fault profiles under the frozen execution protocol?
\end{description}

The evidence is intentionally layered. RQ1 provides bounded property
verification, RQ2 evaluates sensitivity to predefined scientific mutants, RQ3
evaluates bounded implementation-model trace conformance without claiming
refinement or behavioral equivalence, and RQ4 characterizes verification cost
within each tool without implying an asymptotic law or absolute
TLC-versus-Alloy superiority.

The definitive campaign scheduled 1,272 tasks. RQ1 included 420 measured
executions, of which 350 completed and 70 TLC-large executions were censored by
timeout, with no property violations observed among completed configurations.
RQ2 detected all 10 predefined scientific mutants, yielding a mutation score of
1.0 for the evaluated catalogue. RQ3 comprised 600 executions covering valid
and deliberately corrupted traces, valid traces were accepted and corrupted
traces rejected with the expected diagnostics. RQ4 showed increasing
verification cost across completed Alloy profiles, while all measured
TLC-large executions reached timeout. These observations support bounded and
tool-specific conclusions only.

Reproduction used a separate clone, execution workspace, and operating-system
user on the same native Linux host. The successful attempt completed 10/10
predefined gates and reproduced 32/32 expected analysis-artifact hashes. This
demonstrates process separation and deterministic reconstruction of preserved
analytical outputs under the documented environment, not hardware independence,
a complete rerun of all 1,272 tasks, or scientific correctness from hash
agreement alone.

The remainder of the paper presents related work, the protocol and formal
models, methodology and experimental design, results and interpretation,
threats to validity, and reproducibility evidence.
